\documentclass[11pt,a4paper]{article}
\usepackage[top=1.6cm,bottom=1.4cm,left=1.3cm,right=1.3cm,headsep=8pt,footskip=16pt]{geometry}
\usepackage{enumitem}
\usepackage{amssymb}
\usepackage{array}
\usepackage[table]{xcolor}
\usepackage{fancyhdr}
\usepackage{multicol}
\usepackage{polyglossia}
\setdefaultlanguage[calendar=gregorian, numerals=maghrib]{arabic}
\setotherlanguage{english}
\IfFontExistsTF{Amiri}{%
  \newfontfamily\arabicfont[Script=Arabic, Scale=1.08]{Amiri}
}{%
  \newfontfamily\arabicfont[Script=Arabic, Scale=1.08, Path=../fonts/, Extension=.ttf,
    UprightFont=*-Regular, BoldFont=*-Bold, ItalicFont=*-Italic,
    BoldItalicFont=*-BoldItalic]{Amiri}
}
\IfFontExistsTF{AlBattar}{%
  \newfontfamily\khatt[Script=Arabic, Scale=1.15]{AlBattar}
}{%
  \newfontfamily\khatt[Script=Arabic, Scale=1.15, Path=../fonts/,
    BoldFont=ae_AlBattar.ttf]{ae_AlBattar.ttf}
}
\renewcommand{\headrulewidth}{0.8pt}
\renewcommand{\footrulewidth}{0.8pt}
\fancyhf{}
\rhead{\khatt المدرسة العليا للأساتذة ببوزريعة}
\lhead{\khatt ملحقة بومرداس}
\lfoot{\khatt الأستاذ: د. بن عبد الرزاق يوسف}
\cfoot{\thepage/2}
\pagestyle{fancy}
\parindent 0pt
\setlength{\columnsep}{0.6cm}
\setlength{\columnseprule}{0.4pt}
\setlist[enumerate,1]{label=\textbf{\arabic*.}, leftmargin=*, itemsep=5pt, topsep=2pt}
\begin{document}

\begin{center}
{\Large\khatt\textbf{الفرض الثالث في الذكاء الاصطناعي}}\\[2pt]
{\normalsize أسئلة متعددة الاختيارات ــ الفصول من 1 إلى 6 ــ 15 سؤالًا}
\end{center}
\vspace{-4pt}\rule{\textwidth}{0.6pt}\\
{\small\textbf{التعليمة:} لكل سؤال إجابة صحيحة واحدة فقط من بين خمسة اختيارات.
اكتب \textbf{حرف} الإجابة الصحيحة في جدول الإجابات أسفل الصفحة الثانية.
لكل سؤال نقطة واحدة (المجموع 15 نقطة).}
\vspace{-2pt}
\begin{multicols}{2}

\noindent\begin{minipage}{\linewidth}
\hangindent=0.7cm\hangafter=1\makebox[0.7cm][r]{\textbf{1.}}\textbf{أي أداة من الأدوات الآتية تطوّرها شركة \textenglish{Google}؟}
\begin{itemize}[leftmargin=0.95cm, labelsep=4pt, itemsep=0pt, topsep=2pt, parsep=0pt]
  \item[\textbf{أ)}] \textenglish{DeepSeek}
  \item[\textbf{ب)}] \textenglish{ChatGPT}
  \item[\textbf{ج)}] \textenglish{Claude}
  \item[\textbf{د)}] \textenglish{Gemini}
  \item[\textbf{هـ)}] \textenglish{Qwen}
\end{itemize}
\end{minipage}\par\vspace{7pt plus 4pt}
\noindent\begin{minipage}{\linewidth}
\hangindent=0.7cm\hangafter=1\makebox[0.7cm][r]{\textbf{2.}}\textbf{عبارة «التلاميذ في السنة الثانية متوسط» تمثّل عنصر:}
\begin{itemize}[leftmargin=0.95cm, labelsep=4pt, itemsep=0pt, topsep=2pt, parsep=0pt]
  \item[\textbf{أ)}] المهمة
  \item[\textbf{ب)}] القيود
  \item[\textbf{ج)}] السياق
  \item[\textbf{د)}] التنسيق
  \item[\textbf{هـ)}] الدور
\end{itemize}
\end{minipage}\par\vspace{7pt plus 4pt}
\noindent\begin{minipage}{\linewidth}
\hangindent=0.7cm\hangafter=1\makebox[0.7cm][r]{\textbf{3.}}\textbf{ما التلخيص؟}
\begin{itemize}[leftmargin=0.95cm, labelsep=4pt, itemsep=0pt, topsep=2pt, parsep=0pt]
  \item[\textbf{أ)}] إضافة أفكار جديدة إلى النص لجعله أغنى وأطول
  \item[\textbf{ب)}] ترجمة النص كاملًا إلى لغة أخرى دون حذف شيء
  \item[\textbf{ج)}] نسخ النص كاملًا كما هو مع تغيير العنوان فقط
  \item[\textbf{د)}] حذف الفقرة الأولى والأخيرة من النص فقط
  \item[\textbf{هـ)}] كتابة الأفكار الأساسية بأسلوبنا دون تغيير المعنى
\end{itemize}
\end{minipage}\par\vspace{7pt plus 4pt}
\noindent\begin{minipage}{\linewidth}
\hangindent=0.7cm\hangafter=1\makebox[0.7cm][r]{\textbf{4.}}\textbf{أعطينا الطلب نفسه لأداتين مختلفتين. ماذا نتوقع؟}
\begin{itemize}[leftmargin=0.95cm, labelsep=4pt, itemsep=0pt, topsep=2pt, parsep=0pt]
  \item[\textbf{أ)}] رفض إحدى الأداتين للطلب في كل مرة
  \item[\textbf{ب)}] إجابتين مختلفتين في الوضوح والطول والتنظيم
  \item[\textbf{ج)}] إجابة صحيحة من الأولى وخاطئة من الثانية
  \item[\textbf{د)}] لا إجابة، لأن الطلب استُعمل من قبل
  \item[\textbf{هـ)}] إجابتين متطابقتين حرفيًا في كل كلمة
\end{itemize}
\end{minipage}\par\vspace{7pt plus 4pt}
\noindent\begin{minipage}{\linewidth}
\hangindent=0.7cm\hangafter=1\makebox[0.7cm][r]{\textbf{5.}}\textbf{لدينا 1000 مثال نقسمها بنسبة 80٪ للتدريب و20٪ للاختبار. كم عدد أمثلة الاختبار؟}
\begin{itemize}[leftmargin=0.95cm, labelsep=4pt, itemsep=0pt, topsep=2pt, parsep=0pt]
  \item[\textbf{أ)}] 500
  \item[\textbf{ب)}] 100
  \item[\textbf{ج)}] 800
  \item[\textbf{د)}] 200
  \item[\textbf{هـ)}] 250
\end{itemize}
\end{minipage}\par\vspace{7pt plus 4pt}
\noindent\begin{minipage}{\linewidth}
\hangindent=0.7cm\hangafter=1\makebox[0.7cm][r]{\textbf{6.}}\textbf{في الأمر: «تصرّف كأستاذ جامعي في البلاغة. اشرح مفهوم التشبيه لطالب في السنة الأولى...»، عبارة «لطالب في السنة الأولى» تمثّل:}
\begin{itemize}[leftmargin=0.95cm, labelsep=4pt, itemsep=0pt, topsep=2pt, parsep=0pt]
  \item[\textbf{أ)}] الدور
  \item[\textbf{ب)}] المهمة
  \item[\textbf{ج)}] القيود
  \item[\textbf{د)}] التنسيق
  \item[\textbf{هـ)}] السياق
\end{itemize}
\end{minipage}\par\vspace{7pt plus 4pt}
\noindent\begin{minipage}{\linewidth}
\hangindent=0.7cm\hangafter=1\makebox[0.7cm][r]{\textbf{7.}}\textbf{في الدرجة الثانية من بناء الأمر خطوة بخطوة، ما العنصر الذي أُضيف إلى الأمر العام؟}
\begin{itemize}[leftmargin=0.95cm, labelsep=4pt, itemsep=0pt, topsep=2pt, parsep=0pt]
  \item[\textbf{أ)}] القيود والتنسيق
  \item[\textbf{ب)}] الدور
  \item[\textbf{ج)}] المهمة والقيود
  \item[\textbf{د)}] التنسيق
  \item[\textbf{هـ)}] السياق
\end{itemize}
\end{minipage}\par\vspace{7pt plus 4pt}
\noindent\begin{minipage}{\linewidth}
\hangindent=0.7cm\hangafter=1\makebox[0.7cm][r]{\textbf{8.}}\textbf{أي جملة إفصاح مناسبة؟}
\begin{itemize}[leftmargin=0.95cm, labelsep=4pt, itemsep=0pt, topsep=2pt, parsep=0pt]
  \item[\textbf{أ)}] «قد أكون استعملت شيئًا ما أثناء كتابة هذا البحث»
  \item[\textbf{ب)}] «استعملتُ \textenglish{ChatGPT} لتنظيم الأفكار، والصياغة من عملي»
  \item[\textbf{ج)}] «هذا البحث من عمل الأداة كله، وأنا قدّمته فقط»
  \item[\textbf{د)}] «لا أذكر مصادري ولا أدواتي لأنها كثيرة جدًا»
  \item[\textbf{هـ)}] «كتبت الأداة هذا البحث، ولا أتحمل أي خطأ فيه»
\end{itemize}
\end{minipage}\par\vspace{7pt plus 4pt}
\noindent\begin{minipage}{\linewidth}
\hangindent=0.7cm\hangafter=1\makebox[0.7cm][r]{\textbf{9.}}\textbf{أي مما يلي سؤال بحث محدد؟}
\begin{itemize}[leftmargin=0.95cm, labelsep=4pt, itemsep=0pt, topsep=2pt, parsep=0pt]
  \item[\textbf{أ)}] «التعليم والمدرسة والتكنولوجيا الحديثة»
  \item[\textbf{ب)}] «الذكاء الاصطناعي والتعليم في العالم الحديث»
  \item[\textbf{ج)}] «كل شيء عن النحو العربي وتاريخه الطويل»
  \item[\textbf{د)}] «كيف يحسّن الذكاء الاصطناعي التعبير الكتابي؟»
  \item[\textbf{هـ)}] «الأدب العربي القديم والحديث والمعاصر»
\end{itemize}
\end{minipage}\par\vspace{7pt plus 4pt}
\noindent\begin{minipage}{\linewidth}
\hangindent=0.7cm\hangafter=1\makebox[0.7cm][r]{\textbf{10.}}\textbf{في المخطط المفاهيمي لدرس «الجملة الاسمية»، المفهومان المرتبطان مباشرة بالجملة الاسمية هما:}
\begin{itemize}[leftmargin=0.95cm, labelsep=4pt, itemsep=0pt, topsep=2pt, parsep=0pt]
  \item[\textbf{أ)}] الفعل والفاعل
  \item[\textbf{ب)}] المبتدأ والخبر
  \item[\textbf{ج)}] الحال والتمييز
  \item[\textbf{د)}] الجار والمجرور
  \item[\textbf{هـ)}] النعت والمنعوت
\end{itemize}
\end{minipage}\par\vspace{7pt plus 4pt}
\noindent\begin{minipage}{\linewidth}
\hangindent=0.7cm\hangafter=1\makebox[0.7cm][r]{\textbf{11.}}\textbf{قال طالب: «\textenglish{ChatGPT} هو الذكاء الاصطناعي كله». ما التصحيح الأدق؟}
\begin{itemize}[leftmargin=0.95cm, labelsep=4pt, itemsep=0pt, topsep=2pt, parsep=0pt]
  \item[\textbf{أ)}] الذكاء الاصطناعي نوع من أنواع \textenglish{ChatGPT} وجزء منه
  \item[\textbf{ب)}] الذكاء الاصطناعي التوليدي أوسع من الذكاء الاصطناعي
  \item[\textbf{ج)}] \textenglish{ChatGPT} أداة توليدية، والتوليدي نوع من الذكاء الاصطناعي
  \item[\textbf{د)}] \textenglish{ChatGPT} ليس ذكاءً اصطناعيًا، بل محرك بحث تقليدي
  \item[\textbf{هـ)}] كلامه صحيح، لأن \textenglish{ChatGPT} هو الأداة الوحيدة الموجودة
\end{itemize}
\end{minipage}\par\vspace{7pt plus 4pt}
\noindent\begin{minipage}{\linewidth}
\hangindent=0.7cm\hangafter=1\makebox[0.7cm][r]{\textbf{12.}}\textbf{أخذ فريق 20٪ من البيانات للاختبار، لكنه استعمل الأمثلة نفسها في التدريب أيضًا. ما الخلل؟}
\begin{itemize}[leftmargin=0.95cm, labelsep=4pt, itemsep=0pt, topsep=2pt, parsep=0pt]
  \item[\textbf{أ)}] لم يعد الاختبار بأمثلة جديدة، فلا يقيس الفهم
  \item[\textbf{ب)}] الخلل أن البيانات صور وليست نصوصًا مكتوبة
  \item[\textbf{ج)}] الخلل أن نسبة الاختبار صغيرة جدًا وغير كافية
  \item[\textbf{د)}] الخلل في عدد مراحل التعلم التي اتبعها الفريق
  \item[\textbf{هـ)}] لا خلل في ذلك، لأن نسبة 20٪ نسبة محترمة
\end{itemize}
\end{minipage}\par\vspace{7pt plus 4pt}
\noindent\begin{minipage}{\linewidth}
\hangindent=0.7cm\hangafter=1\makebox[0.7cm][r]{\textbf{13.}}\textbf{أي جملة تحقق الإفصاح والشفافية معًا؟}
\begin{itemize}[leftmargin=0.95cm, labelsep=4pt, itemsep=0pt, topsep=2pt, parsep=0pt]
  \item[\textbf{أ)}] «استعملتُ أداة ما في البحث، ولا أذكر في أي جزء استعملتها»
  \item[\textbf{ب)}] «الأداة كتبت كل شيء في هذا البحث، وأنا اكتفيت بتقديمه»
  \item[\textbf{ج)}] «استعملتُ \textenglish{Gemini} لاقتراح خطة الفصل الأول، وتحققت من المعلومات»
  \item[\textbf{د)}] «لم أستعمل أي أداة ذكاء اصطناعي في أي جزء من هذا البحث»
  \item[\textbf{هـ)}] «استعملتُ الذكاء الاصطناعي في هذا البحث، وهذا كل ما أذكره»
\end{itemize}
\end{minipage}\par\vspace{7pt plus 4pt}
\noindent\begin{minipage}{\linewidth}
\hangindent=0.7cm\hangafter=1\makebox[0.7cm][r]{\textbf{14.}}\textbf{أضافت أداة إلى ملخص نصٍّ إحصائيةً غير موجودة في النص الأصلي. أي سؤال مراجعة يكشف ذلك؟}
\begin{itemize}[leftmargin=0.95cm, labelsep=4pt, itemsep=0pt, topsep=2pt, parsep=0pt]
  \item[\textbf{أ)}] هل أضافت معلومة غير موجودة؟
  \item[\textbf{ب)}] هل الملخص قصير بما يكفي؟
  \item[\textbf{ج)}] هل حذفت فكرة مهمة من النص؟
  \item[\textbf{د)}] هل الملخص منسق في فقرات؟
  \item[\textbf{هـ)}] هل الملخص مكتوب بالفصحى؟
\end{itemize}
\end{minipage}\par\vspace{7pt plus 4pt}
\noindent\begin{minipage}{\linewidth}
\hangindent=0.7cm\hangafter=1\makebox[0.7cm][r]{\textbf{15.}}\textbf{شرحت أداة صورة مخطط بيانيًا شرحًا مقنعًا. لماذا لا يكفي ذلك لاستعماله في القسم؟}
\begin{itemize}[leftmargin=0.95cm, labelsep=4pt, itemsep=0pt, topsep=2pt, parsep=0pt]
  \item[\textbf{أ)}] قد تخطئ في الأرقام أو العلاقات، فنقارن بالأصل
  \item[\textbf{ب)}] لأن كل شرح مقنع خاطئ بالضرورة ولا يُعتمد
  \item[\textbf{ج)}] لأن الشرح يجب أن يكون صوتيًا فقط لا مكتوبًا
  \item[\textbf{د)}] لأن الأداة لا تستطيع قراءة المخططات أبدًا
  \item[\textbf{هـ)}] لأن المخططات البيانية لا تُستعمل في القسم
\end{itemize}
\end{minipage}\par\vspace{7pt plus 4pt}
\end{multicols}
\vfill
\begin{center}
{\khatt\textbf{جدول الإجابات}}\\[3pt]
\renewcommand{\arraystretch}{1.6}
\begin{tabular}{|>{\centering\arraybackslash}m{1.9cm}|*{15}{>{\centering\arraybackslash}m{0.62cm}|}}
\hline
\textbf{رقم السؤال} & \textbf{1} & \textbf{2} & \textbf{3} & \textbf{4} & \textbf{5} & \textbf{6} & \textbf{7} & \textbf{8} & \textbf{9} & \textbf{10} & \textbf{11} & \textbf{12} & \textbf{13} & \textbf{14} & \textbf{15} \\ \hline
\textbf{الإجابة} &  &  &  &  &  &  &  &  &  &  &  &  &  &  &  \\[8pt] \hline
\end{tabular}

\vspace{8pt}
\renewcommand{\arraystretch}{2.2}
\begin{tabular}{|>{\centering\arraybackslash}m{7.5cm}|>{\centering\arraybackslash}m{3.2cm}|>{\centering\arraybackslash}m{3.2cm}|}
\hline
\textbf{الاسم واللقب} & \textbf{الفوج} & \textbf{العلامة} \\ \hline
 & & \hfill /15 \\ \hline
\end{tabular}
\end{center}
\end{document}
